% TK2 stationary densities, mean composition and reduction error.
% Data are in data/tk2-correction/.
\documentclass[10pt, border = 2pt]{standalone}
\usepackage{amsmath, amssymb}
\usepackage[dvipsnames]{xcolor}
\usepackage{pgfplots}
\pgfplotsset{compat = 1.18}
\usetikzlibrary{arrows.meta, calc}
% Shared colours, typography and plot styles for the manuscript figures.
\providecommand{\figtick}{\footnotesize}
\providecommand{\figlabel}{\small}
\providecommand{\figheading}{\normalsize}

\definecolor{fastgrey}{HTML}{6F6F6F}
\definecolor{leading}{HTML}{0072B2}
\definecolor{firstorder}{HTML}{A64D79}
\definecolor{stateorange}{HTML}{E69F00}

\tikzset{
    sublabel/.style = {
        font = \figlabel,
        inner sep = 0pt,
    },
    keyline/.style = {
        line width = 0.8pt,
    },
    keytext/.style = {
        font = \figtick,
        text = black!75,
        anchor = west,
        inner sep = 2pt,
    },
    manifold/.style = {
        draw = leading,
        line width = 1.4pt,
        line cap = round,
    },
    fibre/.style = {
        draw = black!18,
        line width = 0.4pt,
    },
    driftbox/.style = {
        draw = black!30,
        line width = 0.5pt,
        rounded corners = 3pt,
    },
    noisebox/.style = {
        draw = leading!70,
        line width = 0.5pt,
        rounded corners = 3pt,
    },
    slowarrow/.style = {
        -{Latex[length = 5pt, width = 4.5pt]},
        draw = leading,
        line width = 1.3pt,
        line cap = round,
    },
    fastarrow/.style = {
        ->,
        draw = fastgrey,
        line width = 0.9pt,
        line cap = round,
        shorten < = 2.5pt,
        shorten > = 2.5pt,
    },
    noisyarrow/.style = {
        fastarrow,
        decorate,
        decoration = {
            snake,
            segment length = 6pt,
            amplitude = 0.8pt,
            pre length = 2pt,
            post length = 3pt,
        },
    },
    state/.style = {
        circle,
        fill = stateorange,
        draw = white,
        line width = 0.6pt,
        minimum size = 0.2cm,
        inner sep = 0pt,
    },
    zero/.style = {
        state,
        fill = leading,
    },
    fixedpt/.style = {
        circle,
        fill = black,
        inner sep = 0pt,
        minimum size = 3.5pt,
    },
}

\ifdefined\pgfplotsset
\pgfplotsset{
    paperaxis/.style = {
        axis line style = {
            black!45,
            line width = 0.4pt,
        },
        tick style = {black!45},
        tick label style = {
            font = \figtick,
            text = black!70,
        },
        label style = {
            font = \figlabel,
            text = black!80,
        },
        axis x line* = bottom,
        axis y line* = left,
        scale only axis,
        axis on top,
    },
    legend image code/.code = {
        \draw[
            mark repeat = 2,
            mark phase = 2,
            #1,
        ] plot coordinates {
            (0cm, 0cm)
            (0.15cm, 0cm)
            (0.3cm, 0cm)
        };
    },
    exactlegend/.style = {
        legend image code/.code = {
            \fill[black!16]
                (0cm, -0.08cm) rectangle (0.3cm, 0.08cm);
            \filldraw[
                fill = black!35,
                draw = black!75,
                line width = 0.3pt,
            ] (0.15cm, 0cm) circle (1.5pt);
        },
    },
    shade/.style = {
        draw = none,
        fill = leading!8,
        const plot,
        forget plot,
    },
    exactpath/.style = {
        black!45,
        line width = 0.5pt,
    },
    exactfill/.style = {
        draw = none,
        fill = black!16,
        const plot,
    },
    exactmark/.style = {
        only marks,
        mark = *,
        mark size = 1.5pt,
        mark options = {
            fill = black!35,
            draw = black!75,
            line width = 0.3pt,
        },
    },
    exactsquare/.style = {
        only marks,
        mark = square*,
        mark size = 1.4pt,
        mark options = {
            fill = black!35,
            draw = black!75,
            line width = 0.3pt,
        },
    },
    leadingline/.style = {
        leading,
        line width = 1.0pt,
    },
    correctedline/.style = {
        firstorder,
        line width = 1.4pt,
    },
}
\fi


\def\data{data/tk2-correction}
\input{\data/params}

\pgfplotsset{
    density/.style = {
        paperaxis,
        width = 5.6cm,
        height = 3.4cm,
        xmin = \sAxisMin,
        xmax = \sAxisMax,
        ymin = 0,
        xtick = {0.5, 1, 1.5},
        ylabel = {$\tilde\rho(s)$},
    },
    epsaxis/.style = {
        xmin = \epsAxisMin,
        xmax = \epsAxisMax,
        xtick = {0.002, 0.005, 0.01, 0.02, 0.05, 0.1, 0.2},
        xticklabels = {0.002, 0.005, 0.01, 0.02, 0.05, 0.1, 0.2},
        xlabel = {$\varepsilon$},
    },
    erroraxis/.style = {
        paperaxis,
        width = 5.75cm,
        height = 3.05cm,
        ylabel = {$W_1$ error},
        yminorticks = true,
    },
    errline/.style = {
        mark = x,
        mark size = 2.3pt,
        mark options = {solid, line width = 0.9pt},
    },
}

\begin{document}
\begin{tikzpicture}
    % (a) Stationary densities at two epsilon values
    \begin{axis}[
        density,
        name = densa,
        ymax = \densityMaxA,
        xticklabels = {},
    ]
        \addplot[
            exactfill,
        ] table[
            x = s,
            y = density,
            col sep = comma,
        ] {\data/exact-a.csv} \closedcycle;
        \addplot[
            correctedline,
        ] table[
            x = s,
            y = corrected,
            col sep = comma,
        ] {\data/density-a.csv};
        \addplot[
            leadingline,
        ] table[
            x = s,
            y = leading,
            col sep = comma,
        ] {\data/density-a.csv};
        \node[
            font = \figtick,
            anchor = north west,
        ] at (rel axis cs: 0.02, 0.98) {$\varepsilon = \epsRowA$};
    \end{axis}

    \begin{axis}[
        density,
        name = densb,
        at = {($(densa.south west) + (0, -0.25cm)$)},
        anchor = north west,
        ymax = \densityMaxB,
        xlabel = {total $s$},
    ]
        \addplot[
            exactfill,
        ] table[
            x = s,
            y = density,
            col sep = comma,
        ] {\data/exact-b.csv} \closedcycle;
        \addplot[
            correctedline,
        ] table[
            x = s,
            y = corrected,
            col sep = comma,
        ] {\data/density-b.csv};
        \addplot[
            leadingline,
        ] table[
            x = s,
            y = leading,
            col sep = comma,
        ] {\data/density-b.csv};
        \node[
            font = \figtick,
            anchor = north west,
        ] at (rel axis cs: 0.02, 0.98) {$\varepsilon = \epsRowB$};
    \end{axis}
    \node[
        sublabel,
        anchor = north west,
    ] at ($(densa.north west) + (-1.05, 0)$) {(a)};

    % (b) Mean relative composition
    \begin{axis}[
        paperaxis,
        epsaxis,
        name = mean,
        at = {($(densa.north east) + (1.55cm, 0)$)},
        anchor = north west,
        width = 5.75cm,
        height = 2.0cm,
        xmode = log,
        ymin = \meanYMin,
        ymax = \meanYMax,
        ylabel = {$\langle m\rangle$},
        yticklabel style = {
            /pgf/number format/fixed,
            /pgf/number format/precision = 2,
        },
        legend style = {
            at = {(0.03, 0.04)},
            anchor = south west,
            draw = none,
            fill = white,
            fill opacity = 0.85,
            text opacity = 1,
            font = \figtick,
            legend columns = 1,
            row sep = -1pt,
        },
        legend cell align = left,
    ]
        \addlegendimage{exactlegend}
        \addlegendentry{exact}
        \addlegendimage{leadingline}
        \addlegendentry{$O(1)$}
        \addlegendimage{correctedline}
        \addlegendentry{$O(\varepsilon)$}

        \draw[
            black!22,
            line width = 0.4pt,
        ] (axis cs: \epsRowA, \meanYMin)
            -- (axis cs: \epsRowA, \meanYMax);
        \draw[
            black!22,
            line width = 0.4pt,
        ] (axis cs: \epsRowB, \meanYMin)
            -- (axis cs: \epsRowB, \meanYMax);

        \addplot[
            correctedline,
            forget plot,
        ] table[
            x = epsilon,
            y = y,
            col sep = comma,
        ] {\data/mean-corrected.csv};
        \addplot[
            leadingline,
            domain = \epsAxisMin:\epsAxisMax,
            samples = 2,
            forget plot,
        ] {\leadingMeanY};
        \addplot[
            exactmark,
            forget plot,
        ] table[
            x = epsilon,
            y = y,
            col sep = comma,
        ] {\data/mean-exact.csv};
    \end{axis}
    \node[
        sublabel,
        anchor = north west,
    ] at ($(mean.north west) + (-1.05, 0)$) {(b)};

    % (c) Wasserstein error at two volumes
    \path let
        \p1 = (mean.south),
        \p2 = (densb.south)
    in \pgfextra{
        \pgfmathsetlengthmacro\errHeight{\y1 - \y2 - 1.15cm}
        \global\let\errHeight\errHeight
    };
    \begin{axis}[
        erroraxis,
        epsaxis,
        name = errc,
        at = {($(mean.south west) + (0, -1.15cm)$)},
        anchor = north west,
        height = \errHeight,
        xmode = log,
        ymode = log,
        ymin = \epsErrorMin,
        ymax = \epsErrorMax,
        legend style = {
            at = {(0.03, 0.97)},
            anchor = north west,
            draw = none,
            fill = none,
            font = \figtick,
            legend columns = 1,
        },
        legend cell align = left,
    ]
        \addlegendimage{black!70, dashed, line width = 0.8pt}
        \addlegendentry{$V = \errorVolumeOne$}
        \addlegendimage{black!70, line width = 0.8pt}
        \addlegendentry{$V = \errorVolumeTwo$}

        % Reference lines have slopes one and two.
        \addplot[
            black!40,
            line width = 0.5pt,
            domain = \epsGuideX/10:\epsGuideX,
            samples = 2,
            forget plot,
        ] {\leadingGuideY*(x/\epsGuideX)};
        \addplot[
            black!40,
            line width = 0.5pt,
            domain = \epsGuideX/10:\epsGuideX,
            samples = 2,
            forget plot,
        ] {\correctedGuideY*(x/\epsGuideX)^2};

        \node[
            font = \figtick,
            text = black!55,
            anchor = north west,
        ] at (
            axis cs: \epsGuideX/3.16,
            \leadingGuideY/3.16
        ) {$\varepsilon$};
        \node[
            font = \figtick,
            text = black!55,
            anchor = north west,
        ] at (
            axis cs: \epsGuideX/3.16,
            \correctedGuideY/10
        ) {$\varepsilon^2$};

        \addplot[
            leadingline,
            errline,
            dashed,
            forget plot,
        ] table[
            x = epsilon,
            y = leadingone,
            col sep = comma,
        ] {\data/errors-epsilon.csv};
        \addplot[
            correctedline,
            errline,
            dashed,
            forget plot,
        ] table[
            x = epsilon,
            y = correctedone,
            col sep = comma,
        ] {\data/errors-epsilon.csv};
        \addplot[
            leadingline,
            errline,
            forget plot,
        ] table[
            x = epsilon,
            y = leadingtwo,
            col sep = comma,
        ] {\data/errors-epsilon.csv};
        \addplot[
            correctedline,
            errline,
            forget plot,
        ] table[
            x = epsilon,
            y = correctedtwo,
            col sep = comma,
        ] {\data/errors-epsilon.csv};
    \end{axis}
    \node[
        sublabel,
        anchor = north west,
    ] at ($(errc.north west) + (-1.05, 0)$) {(c)};
\end{tikzpicture}
\end{document}
